% Converted from README.md
\subsection{Laboratorijski primeri}\label{dijagnoziranje-primeri}

Primeri su preuzeti iz \href{https://github.com/brendangregg/perf-labs}{perf-labs} repozitorijuma Brendana Gregg-a, korišćenog tokom \href{https://youtu.be/FJW8nGV4jxY}{prezentacije o alatima za merenje performansi Linux kernel-a}. Svaki program namerno pravi jedan simptom. Cilj vežbe nije odmah otvoriti izvorni kod, već prvo posmatrati proces kao crnu kutiju i zatim potvrditi uzrok.

Pre pokretanja primera treba ih prevesti:

\begin{minted}{bash}
$ make all
\end{minted}

Pokretati jedan primer po jedan, u posebnom terminalu. Posle analize zaustaviti proces pomoću \texttt{Ctrl+C} ili komandom:

\begin{minted}{bash}
$ pkill lab001
\end{minted}

Tok dijagnoze za svaki primer je isti:

\begin{enumerate}
\tightlist
\item
  proveriti opšte stanje sistema;
\item
  pronaći proces koji pravi simptom;
\item
  odrediti resurs koji je usko grlo;
\item
  potvrditi uzrok alatom koji posmatra bliži nivo sistema.
\end{enumerate}

\subsubsection{Pregled primera}\label{dijagnoziranje-labs-pregled}

\begin{longtable}[]{@{}lll@{}}
\toprule
Primer & Simptom & Prvi alati \\
\midrule
\endhead
\texttt{lab001} & CPU petlja u jednoj funkciji & \texttt{top}, \texttt{pidstat}, \texttt{strace} \\
\texttt{lab002} & proces čeka konekciju & \texttt{ss}, \texttt{strace}, \texttt{tcpdump} \\
\texttt{lab003} & veliki broj sistemskih poziva & \texttt{vmstat}, \texttt{pidstat}, \texttt{strace} \\
\texttt{lab004} & kratkotrajni child procesi & \texttt{pidstat}, \texttt{ps}, \texttt{strace -f} \\
\texttt{lab005} & sinhroni upisi na disk & \texttt{vmstat}, \texttt{iostat}, \texttt{pidstat -d} \\
\texttt{lab006} & CPU rad raspodeljen kroz više funkcija & \texttt{top}, \texttt{mpstat}, \texttt{pidstat} \\
\texttt{lab007} & proces uglavnom spava & \texttt{top}, \texttt{pidstat -w}, \texttt{strace} \\
\texttt{lab008} & upisi na disk sa povremenim kašnjenjem & \texttt{iostat}, \texttt{pidstat -d}, \texttt{strace -T} \\
\texttt{lab013} & iscrpljivanje steka rekurzijom & \texttt{ulimit}, \texttt{dmesg}, \texttt{sar -r} \\
\bottomrule
\end{longtable}

\subsubsection{\texorpdfstring{\texttt{lab001}: CPU petlja}{lab001: CPU petlja}}\label{lab001-cpu-petlja}

Ovaj program zauzima CPU bez sistemskih poziva i bez I/O aktivnosti. Kod je mali, ali je koristan za vežbanje razlike između procesa koji radi u korisničkom prostoru i procesa koji čeka kernel ili uređaj.

\codefile{lab001.c}{c}

Pokrenuti program i posmatrati ga iz drugog terminala:

\begin{minted}{bash}
$ ./lab001 &
$ top -p $(pgrep lab001)
$ pidstat -p $(pgrep lab001) 1
$ strace -p $(pgrep lab001)
\end{minted}

Očekivano je da \texttt{top} i \texttt{pidstat} pokažu visoku CPU potrošnju. \texttt{strace} neće prikazivati stalne sistemske pozive, jer se program vrti u funkciji \texttt{func\_c}. Time se problem svodi na CPU-bound kod u korisničkom prostoru. Za tačnu lokaciju vrućeg steka koriste se debager ili profajler.

\subsubsection{\texorpdfstring{\texttt{lab002}: čekanje na konekciju}{lab002: čekanje na konekciju}}\label{lab002-cekanje-na-konekciju}

Ovaj program otvara TCP port i zatim čeka klijentsku konekciju. Simptom je proces koji postoji, ali troši malo CPU vremena i deluje neaktivan.

\codefile{lab002.c}{c}

Prvo proveriti da li proces sluša na očekivanom portu:

\begin{minted}{bash}
$ ./lab002 &
$ ss -ltnp 'sport = :12340'
$ lsof -iTCP:12340 -sTCP:LISTEN -P -n
\end{minted}

Zatim proveriti na kom sistemskom pozivu proces čeka:

\begin{minted}{bash}
$ strace -tt -T -p $(pgrep lab002)
\end{minted}

Očekivani trag pokazuje da proces čeka u \texttt{accept}. Kada se klijent poveže, \texttt{accept} se završava, proces čita podatke i šalje odgovor.

\begin{minted}{bash}
$ nc 127.0.0.1 12340
\end{minted}

Ako klijent ne može da se poveže, proveriti port pomoću \texttt{ss}, lokalni firewall i paketni trag pomoću \texttt{tcpdump}. Ako port nije otvoren, problem je pre mrežnog toka: proces nije pokrenut, nije uspeo \texttt{bind}, ili sluša na drugoj adresi.

\subsubsection{\texorpdfstring{\texttt{lab003}: čitanje nula bajtova}{lab003: čitanje nula bajtova}}\label{lab003-citanje-nula-bajtova}

Ovaj program pravi datoteku, otvara je za čitanje i zatim u beskonačnoj petlji poziva \texttt{read} sa veličinom 0. Poziv uspeva, ali ne prenosi podatke. Simptom je neobično velika aktivnost sistemskih poziva bez stvarnog disk I/O rada.

\codefile{lab003.c}{c}

Krenuti od opšteg pregleda:

\begin{minted}{bash}
$ ./lab003 &
$ vmstat 1
$ pidstat -p $(pgrep lab003) 1
$ iostat -x 1
$ sar -n DEV 1
\end{minted}

Ako CPU raste, a disk i mreža ne pokazuju odgovarajući promet, suziti analizu na sistemske pozive:

\begin{minted}{bash}
$ strace -tp $(pgrep lab003) 2>&1 | head -n 20
$ strace -c -p $(pgrep lab003)
\end{minted}

Očekivani nalaz je ponavljanje poziva nalik na \texttt{read(fd, ..., 0) = 0}. To objašnjava zašto proces troši CPU i sistemsko vreme, a ne proizvodi korisno čitanje podataka.

\subsubsection{\texorpdfstring{\texttt{lab004}: kratkotrajni child procesi}{lab004: kratkotrajni child procesi}}\label{lab004-kratkotrajni-procesi}

Ovaj program u petlji pokreće komandu \texttt{cksum} nad datotekom. Problem može lako promaći u \texttt{top}-u, jer child procesi kratko traju.

\codefile{lab004.c}{c}

Početi od procesa roditelja, a zatim pratiti child procese:

\begin{minted}{bash}
$ ./lab004 &
$ pidstat -p $(pgrep lab004) 1
$ ps -ef f | grep -E 'lab004|cksum'
$ strace -f -e trace=process -p $(pgrep lab004)
\end{minted}

Opcija \texttt{-f} je bitna jer prati child procese. Očekivani trag pokazuje ponavljanje kreiranja procesa, izvršavanje \texttt{cksum} i čekanje da child proces završi. Ako se gleda samo roditeljski proces, CPU ili I/O slika može izgledati nepotpuno.

Za proveru disk aktivnosti koristiti:

\begin{minted}{bash}
$ pidstat -d 1
$ iostat -x 1
\end{minted}

\subsubsection{\texorpdfstring{\texttt{lab005}: sinhroni upisi na disk}{lab005: sinhroni upisi na disk}}\label{lab005-sinhroni-upisi}

Ovaj program otvara datoteku sa \texttt{O\_SYNC} i stalno piše u nju. Svaki upis čeka da podaci budu sinhronizovani sa storage slojem. Simptom može biti visok \texttt{iowait}, zauzet disk i spor odziv drugih procesa.

\codefile{lab005.c}{c}

Za ovaj primer je zgodna USE metodologija: proveriti iskorišćenost, zasićenje i greške za CPU, memoriju, disk i mrežu.

\begin{minted}{bash}
$ ./lab005 &
$ htop
$ mpstat -P ALL 1
$ vmstat 1
$ iostat -x 1
$ sar -n DEV 1
\end{minted}

Ako CPU nije potpuno zauzet, nema swap aktivnosti, a \texttt{iostat} pokazuje veliki \texttt{\%util} ili rast \texttt{await}, uzrok je najverovatnije disk. Zatim proveriti procesni I/O i otvorenu datoteku:

\begin{minted}{bash}
$ pidstat -d -p $(pgrep lab005) 1
$ lsof -p $(pgrep lab005) | grep lab005.data
\end{minted}

Bitna veza sa kodom je opcija \texttt{O\_SYNC}. Ona menja cenu svakog \texttt{write} poziva, pa problem nije samo u količini podataka već i u načinu pisanja.

\subsubsection{\texorpdfstring{\texttt{lab006}: više CPU vrućih funkcija}{lab006: više CPU vrućih funkcija}}\label{lab006-cpu-vise-funkcija}

Ovaj program troši CPU u dve funkcije. Za razliku od \texttt{lab001}, vruće vreme nije vezano za jednu jedinu petlju. Primer pokazuje granicu sistemske dijagnoze: alati kao \texttt{top} i \texttt{pidstat} mogu da potvrde CPU problem, ali ne mogu sami da raspodele vreme po funkcijama.

\codefile{lab006.c}{c}

Prvo potvrditi da je problem CPU-bound:

\begin{minted}{bash}
$ ./lab006 &
$ top -p $(pgrep lab006)
$ mpstat -P ALL 1
$ pidstat -p $(pgrep lab006) 1
$ strace -p $(pgrep lab006)
\end{minted}

Ako \texttt{strace} ne prikazuje stalne sistemske pozive, a CPU potrošnja je visoka, proces radi u korisničkom prostoru. Sledeći korak je uzorkovanje steka pomoću profajlera ili analiza procesa u debageru. Ovaj primer zato povezuje poglavlje o dijagnozi sa poglavljima o debagovanju i profajliranju.

\subsubsection{\texorpdfstring{\texttt{lab007}: proces uglavnom spava}{lab007: proces uglavnom spava}}\label{lab007-proces-spava}

Ovaj program periodično spava, pa zatim kratko radi u petlji. Simptom je proces koji je prisutan, ali troši malo CPU vremena. Takav proces može izgledati nebitno u \texttt{top}-u, iako je važan za razumevanje kašnjenja.

\codefile{lab007.c}{c}

Posmatrati CPU vreme i promene konteksta:

\begin{minted}{bash}
$ ./lab007 &
$ top -p $(pgrep lab007)
$ pidstat -w -p $(pgrep lab007) 1
$ strace -tt -T -p $(pgrep lab007)
\end{minted}

Očekivani trag pokazuje periodične pozive za spavanje, najčešće kroz \texttt{nanosleep} ili srodan sistemski poziv. \texttt{pidstat -w} pokazuje promene konteksta, jer proces dobrovoljno prepušta CPU dok spava.

\subsubsection{\texorpdfstring{\texttt{lab008}: disk latencije sa izuzecima}{lab008: disk latencije sa izuzecima}}\label{lab008-disk-latencije}

Ovaj program takođe koristi \texttt{O\_SYNC}, ali povremeno piše veći blok podataka. Simptom nije samo zauzet disk, već i povremeni skokovi latencije pojedinačnih \texttt{write} poziva.

\codefile{lab008.c}{c}

Prvo potvrditi da je usko grlo disk:

\begin{minted}{bash}
$ ./lab008 &
$ iostat -x 1
$ pidstat -d -p $(pgrep lab008) 1
\end{minted}

Zatim izmeriti trajanje sistemskih poziva:

\begin{minted}{bash}
$ strace -tt -T -e trace=write,lseek -p $(pgrep lab008)
\end{minted}

Opcija \texttt{-T} dodaje trajanje svakog sistemskog poziva. Očekivano je da se vide obični upisi i povremeni znatno skuplji upisi. U kodu se to poklapa sa granom koja periodično šalje veći bafer.

\subsubsection{\texorpdfstring{\texttt{lab013}: rekurzija i stek}{lab013: rekurzija i stek}}\label{lab013-rekurzija-stek}

Ovaj program beskonačno rekurzivno poziva istu funkciju. Simptom može izgledati kao memorijski problem, ali uzrok nije heap alokacija već rast steka do ograničenja procesa.

\codefile{lab013.c}{c}

Pošto proces može brzo da se završi, prvo proveriti ograničenje steka i poruke kernela posle pada:

\begin{minted}{bash}
$ ulimit -s
$ ./lab013
$ dmesg -T | tail -n 20
\end{minted}

Ako sistem generiše core dump, sledeći korak je analiza core datoteke u debageru. Ako se želi posmatranje dok proces radi, može se privremeno povećati ograničenje steka u shell-u iz kog se pokreće primer:

\begin{minted}{bash}
$ ulimit -s 65536
$ ./lab013 &
$ pidstat -r -p $(pgrep lab013) 1
\end{minted}

Ovaj primer je dobar podsetnik da globalna memorijska slika ne mora odmah da pokaže uzrok. Lokalna ograničenja procesa, kao što je veličina steka, mogu biti dovoljna da program padne.
